\section{A revised report and the economic history it describes}
\label{sim:financial-source-continuity}

A financial report can change between the two decisions without changing the economic event it describes. Suppose an earlier assertion reports an operating receipt of 7 and a later assertion corrects it to 6. Selecting the available assertion at each date answers an information question. It does not yet determine whether the receipt was always 6, whether an agreement actually changed, or what information the manager received. Those questions must be separated before the selected assertions can support a simulation.

The joint source calculation can now retain a different financial head at each decision. A correction available between the dates selects the old assertion first and the corrected assertion second. With an interval-valued availability time, the supported sequences can be old/old, old/new and new/new. The reverse new/old sequence is impossible under this one-way correction history, even though it would appear in an unconstrained Cartesian product of the two marginal selections. Conflict and retraction retain unresolved content. None of these source selections, by itself, specifies a financial policy.\footnote{The checked financial-source selection companion is retained in \path{data/accounting-master/joint-financial-source-20261027/}. It reproduces 11 protected prior cases and checks 13 new cases, with 341 assertions and 16 alteration controls. Its declared financial-policy binding remains unassessed.}

\subsection{Continuity concerns the complete carried state}

The financing example gives a precise way to distinguish an assertion revision from a change in economic history. Opening cash is $11/2$. The first coupon $I_1$ is accrued and settled; the operating receipt $R$ arrives next; any first-period facility borrowing is repaid; and the second coupon $I_2$ is then accrued and settled. A draw supplies 1 and immediately costs a fee of $1/4$. Let $d_1$ and $d_2$ indicate a draw at the respective decisions. The once-only right and the other payment guards still determine whether an action sequence is permitted.

After the first payment, cash is $11/2-I_1+3d_1/4$. Receipt and repayment change it to $11/2-I_1+R-d_1/4$: repaying the principal removes the first draw's unit of cash, leaving its fee. The second settlement therefore gives
\begin{equation}
 C_{\mathrm{final}}=\frac{11}{2}+R-I_1-I_2-\frac{d_1}{4}+\frac{3d_2}{4}.
 \label{eq:receipt-world-cash}
\end{equation}
The last term includes the second draw's principal because the sequence ends at the second settlement. It does not erase the associated liability. Under receipt 6, fixed coupons 5 and 5 with pay/pay give final cash $3/2$ and no facility fee. Floating coupons 6 and 5 with draw/pay give final cash $1/4$ and a fee of $1/4$. Both examples satisfy the inherited pre- and postconditions. The equation also describes other action sequences algebraically; their guards must be checked separately. All these monetary amounts are supplied synthetic units.

Before the receipt arrives, the complete states under $R=6$ and $R=7$ coincide in this example. That includes the book, contractual memory and recorded flows. Attaching the receipt-6 continuation to the identical pre-receipt prefix therefore reconstructs exactly the complete receipt-6 path. Rejecting every change of a receipt parameter would miss this valid correspondence. The correspondence can support an explicitly assumed correction of the same event; it cannot establish that the revised report is true.

After receipt settlement, the two states differ. Replacing 7 by 6 then removes one unit of cash, changes the income credit and rewrites the recorded operating receipt. Both books balance, but the complete states no longer join. A different contractual right is visible even earlier: the right is carried in opening memory, whether or not it changes the numeric accounts. An actual amendment consequently needs an explicit event and its accounting treatment.\footnote{\path{data/accounting-master/financial-assertion-continuity-20261027/} checks 16 symbolic cases against an independent full-state recurrence: 112 boundary comparisons, 594 assertions and 48 rejected post-receipt joins. It also checks the two viable receipt-6 examples.}

\subsection{An explicit interpretation selects a complete world}

For a supplied same-event correction, the analyst can declare the corrected head, a fixed receipt and a contractual-right family. This declaration supplies a hypothesis about the whole economic history. The selected heads at both dates must lie in the corrected head's receipt-only ancestry, and the right family must remain unchanged. Retaining the superseded assertion as evidence does not impose both 7 and 6 as simultaneous exact facts about one receipt; imposing those incompatible facts would leave no world.

The binding then selects the complete financial rows for the declared world while retaining the dated source pairs, their supports and any unresolved temporal remainder. Receipt 6 with the committed right binds 980 action-region rows; receipt 7 binds 968. These are complete action-region row families, not counts of feasible policies. Their populations agree with the protected reference problems. Unknown interpretation, actual amendment, changed rights, uncovered ancestry, conflict or retraction, inconsistent clocks and incomplete search receive no financial row set.\footnote{The \texttt{SuppliedReceiptCorrectionWorld/1} declaration and its checks are retained in \path{data/accounting-master/financial-source-world-binding-20261027/}: 20 cases, 289 assertions and 25 alteration controls. Twelve cases bind rows and eight remain unresolved or inconsistent. The declaration records its provenance as a supplied synthetic assumption.}

The corrected head may be available to the analyst who specifies the hypothesis without having been available to the earlier manager. Neither that head's identity nor the assumed world's identity is automatically an observable message. An adaptive policy still needs an explicit account of the information available at each decision and must choose the same action whenever that information is the same. The next comparison must therefore preserve the distinction between the analyst's economic hypothesis and the manager's dated observations.
